\documentclass[11pt,a4paper]{article}
\usepackage[utf8]{inputenc}
\usepackage[T1]{fontenc}
\usepackage{lmodern}
\usepackage[margin=1in]{geometry}
\usepackage[dvipsnames]{xcolor}
\usepackage{amsmath,amssymb,amsthm,mathtools}
\usepackage{enumitem}
\usepackage{url}
\usepackage{hyperref}
\usepackage{microtype}
\hypersetup{colorlinks=true,linkcolor=blue,citecolor=teal,urlcolor=blue,
  pdftitle={Sharp symmetry bounds for real algebraic curves},
  pdfauthor={Carlo Perassi}}
\DeclareUrlCommand\lean{\urlstyle{tt}}
\setlist{leftmargin=2em,itemsep=0.25ex,topsep=0.5ex}
\allowdisplaybreaks % chktex 1
\newtheorem{theorem}{Theorem}
\newtheorem{lemma}[theorem]{Lemma}
\theoremstyle{definition}
\newtheorem{remark}[theorem]{Remark}
\newcommand{\C}{\mathbb{C}}
\newcommand{\R}{\mathbb{R}}
\newcommand{\PP}{\mathbb{P}}
\DeclareMathOperator{\Sym}{Sym}
\DeclareMathOperator{\Rea}{Re}
\title{Sharp symmetry bounds for real algebraic curves}
\author{Carlo Perassi}
\date{September 5, 2026 --- revised September 30, 2026}

\begin{document}
\maketitle

\begin{abstract}
Let $C$ be an infinite real plane algebraic curve defined by a polynomial
irreducible over $\C$, of degree $d\ge2$, and suppose that $C$ is not a circle.
Its rotation group has order at most $\max\{d,2d-4\}$, whereas its full
Euclidean symmetry group has order at most $2d$. Both bounds are sharp.
For $d\ge5$ we classify the curves attaining the rotation bound: up to
orientation-preserving similarity they are
$\Rea(z^{d-2}(|z|^2+\alpha))=0$, where $|\alpha|=1$ and
$\alpha\notin\R$. Their spherical closures have a dihedral group of
$4d-8$ holomorphic M\"obius symmetries and no anti-M\"obius symmetry.
The parameter is a complete invariant under holomorphic M\"obius equivalence,
and is determined up to conjugation under full M\"obius equivalence.
An irreducible quintic with a rotation of order six illustrates why a
symmetry of a zero set need not fix its defining polynomial.
Every statement is checked in Lean~4 with Mathlib
(Appendix~\ref{app:lean}).
\end{abstract}

\section{Statements and context}

Throughout, $f\in\R[x,y]$ is irreducible in $\C[x,y]$,
$d=\deg f\ge2$, and $C=\{(x,y)\in\R^2:f(x,y)=0\}$ is infinite and
is not a circle. The degree is the ordinary total degree, not a degree
invariant under M\"obius transformations. Write $\Sym(C)$ for the group
of Euclidean isometries preserving $C$, and $\Sym^+(C)$ for its
orientation-preserving subgroup. The identity is counted in both groups.

\begin{theorem}[Sharp bounds]\label{thm:bounds}
The groups $\Sym(C)$ and $\Sym^+(C)$ are finite. If
$N=|\Sym^+(C)|$, then
\[
 N\le\max\{d,2d-4\},\qquad |\Sym(C)|\le2d.
\]
Both bounds are attained in every degree $d\ge2$.
If $d\ge5$ and $N=2d-4$, then an orientation-preserving similarity
carries $C$ onto
\begin{equation}\label{eq:family}
 C_{m,\alpha}=\{z\in\C:\Rea(z^m(|z|^2+\alpha))=0\},
 \qquad m=d-2,\quad |\alpha|=1,\quad \alpha\notin\R.
\end{equation}
Conversely every curve in~\eqref{eq:family} with $m\ge3$ attains that
rotation bound.
\end{theorem}

The coefficient method behind the proof is classical. Lebmeir and
Richter-Gebert~\cite[Section~4]{LRG}, and Lebmeir~\cite[Theorem~5.4]{Lebmeir},
recover rotational symmetry from the differences of the weights $a-b$
of the nonzero coefficients of $z^a\overline z^{\,b}$. We use this criterion
with irreducibility to obtain degree bounds and their equality cases;
we do not propose a new symmetry-detection method.
Pach and de Zeeuw~\cite[Lemma~2.6]{PZ} use a bound of $4d$ on the number
of isometries in a distinct-distance argument. The bound $2d$ above
sharpens that auxiliary estimate; no improvement of their distance
exponent is claimed.

The distinction between fixing a polynomial and fixing its zero set
matters even for an irreducible curve. The quintic
\begin{equation}\label{eq:quintic}
 f(x,y)=(x^2+y^2)(x^3-3xy^2)-3x^2y+y^3
       =\Rea\bigl(z^3(|z|^2+i)\bigr)
\end{equation}
satisfies $f(e^{\pi i/3}z)=-f(z)$ and has exactly six rotations.
Its irreducibility follows from Lemma~\ref{lem:geometry} below.
Thus it is a counterexample to the degree restriction in~%
\cite[Lemma~9]{ALVv1}, where the direct symmetry group is asserted to
preserve a regular polygon with at most $\deg C$ vertices. This comparison
concerns the cited arXiv version; the journal version has not been checked.

The equality cases have more structure than their Euclidean symmetries
alone reveal. Let $\widehat C_{m,\alpha}$ denote the closure of
$C_{m,\alpha}$ in the Riemann sphere. A full M\"obius equivalence means
either a holomorphic M\"obius map or such a map followed by complex
conjugation, acting on the ambient sphere.

\begin{theorem}[The extremal family]\label{thm:family}
For every integer $m\ge2$ and every nonreal $\alpha$ of modulus one,
$C_{m,\alpha}$ is an infinite geometrically irreducible curve of degree
$m+2$. The full ambient M\"obius symmetry group of $\widehat C_{m,\alpha}$ consists of
\begin{equation}\label{eq:group}
 z\longmapsto cz\quad(c^{2m}=1),\qquad
 z\longmapsto c/z\quad(c^{2m}=\overline\alpha^{\,2}).
\end{equation}
All these maps are holomorphic. The group is dihedral of order $4m$;
the Euclidean symmetry group consists of the $2m$ rotations in the first
set. For fixed $m$, the curves $\widehat C_{m,\alpha}$ and
$\widehat C_{m,\beta}$ are holomorphically M\"obius equivalent if and
only if $\beta=\alpha$, and anti-M\"obius equivalent if and only if
$\beta=\overline\alpha$.
\end{theorem}

The absence of anti-M\"obius symmetries is an ambient statement, not a
claim about the absence of a real structure on the complexified curve.
Neither theorem concerns arbitrary automorphisms of its normalization.

\section{The degree bounds}

We first recall the finiteness argument, including the hypotheses needed
to pass from a real zero set to a polynomial identity.

\begin{lemma}\label{lem:sign}
The group $\Sym(C)$ is finite, its direct subgroup is cyclic, and
every $T\in\Sym(C)$ satisfies $f\circ T=\varepsilon(T)f$ with
$\varepsilon(T)\in\{1,-1\}$. If $T$ is a reflection then
$\varepsilon(T)=1$.
\end{lemma}

\begin{proof}
A nonzero translation preserving $C$ would put infinitely many points
of $C$ on a line: iterate any point of $C$ along the translation vector.
By B\'ezout and irreducibility $C$ would be that line. A nontrivial
glide reflection is excluded by squaring it.
Two rotations with distinct centers cannot both preserve $C$.
Indeed, write them as $R_a(z)=\lambda z+(1-\lambda)a$ and
$R_b(z)=\mu z+(1-\mu)b$, where $\lambda,\mu\ne1$.
Their commutator is a translation, and is nontrivial because
$R_{a}R_{b}$ and $R_{b}R_{a}$ differ by
$(1-\lambda)(1-\mu)(a-b)\ne0$.
Thus all nonidentity rotations have one center $a$.
The images of any point of $C\setminus\{a\}$ under distinct rotations
are distinct points on one circle. B\'ezout bounds their number by $2d$.
If there are no rotations, the direct subgroup is trivial.
In either case it is a finite subgroup of the circle group, hence cyclic.
The full group has index at most two over this subgroup.

The real locus is infinite, hence Zariski dense in the irreducible
complex curve $f=0$. Consequently $f\circ T$ is a nonzero scalar
multiple of $f$, since these polynomials have the same degree.
The scalar is real; iterating the finite-order map $T$ makes it $1$ or
$-1$. If a reflection had scalar $-1$, then $f$ would vanish on its
entire fixed line. That would give a linear factor, contrary to $d\ge2$
and irreducibility.
\end{proof}

Identify $z=x+iy$ and set
\[
 P(X,Y)=f\bigl((X+Y)/2,(X-Y)/(2i)\bigr)
       =\sum_{a+b\le d}p_{ab}X^{a}Y^{b}.
\]
The polynomial $P$ is irreducible and $p_{ba}=\overline{p_{ab}}$.
Translate a rotation center to zero. For a primitive $N$th root
$\zeta$, invariance of $C$ under $z\mapsto\zeta z$ says
\begin{equation}\label{eq:weights}
 P(\zeta X,\zeta^{-1}Y)=\varepsilon P(X,Y),\qquad
 p_{ab}\ne0\ \Longrightarrow\ \zeta^{a-b}=\varepsilon.
\end{equation}

\begin{proof}[Proof of Theorem~\ref{thm:bounds}]
Finiteness follows from Lemma~\ref{lem:sign}. If $N=1$ the bounds are
immediate, so assume $N>1$ and apply~\eqref{eq:weights} to a generator.
If $\varepsilon=1$ and $N>d$, every exponent has $a=b$.
Thus $P=R(XY)$. Irreducibility forces $R$ to be linear, and the infinite
real locus forces $C$ to be a circle, a contradiction. Therefore
$\varepsilon=1$ implies $N\le d$.

Suppose $N>d$ instead. Then $\varepsilon=-1$, $N=2m$ is even, and
$a-b\equiv m\pmod{2m}$. Since $m>d/2$, the only possible weights are
$m$ and $-m$. Hence
\begin{equation}\label{eq:radial}
 P=X^m A(XY)+Y^m\overline A(XY),\qquad
 \deg A\le\left\lfloor\frac{d-m}{2}\right\rfloor,
\end{equation}
where $\overline A$ means coefficientwise conjugation.
If $A$ were constant, $P$ would be a homogeneous binary form of degree
$m\ge2$, hence reducible. Thus $\deg A\ge1$, giving
$m\le d-2$ and $N\le2d-4$.

If there is a reflection $S$, every rotation $R$ is a product of the
reflections $RS$ and $S$. Both fix $f$ by Lemma~\ref{lem:sign}, so
the generator has $\varepsilon=1$. Therefore $N\le d$ and
$|\Sym(C)|=2N\le2d$. Without a reflection,
$|\Sym(C)|=N\le\max\{d,2d-4\}<2d$.

For sharpness of $2d$, use $\Rea(z^d)=1$ for each $d\ge2$.
Its complexification $X^d+Y^d-2$ is irreducible by Eisenstein at
$Y-\eta$ in $\C[Y][X]$, where $\eta^d=2$. The real locus is infinite,
and its $d$ rotations and $d$ reflections give equality.
These examples also attain the rotation bound for $d=2,3$.
For every $d\ge4$, take $m=d-2$ and $\alpha=i$ in
Theorem~\ref{thm:family}, whose proof is independent of the present
theorem. It gives exactly $2d-4$ rotations.

If $d\ge5$ and $N=2d-4>d$, equation~\eqref{eq:radial}
has $m=d-2\ge3$ and $A(s)=a+bs$ with $a,b\ne0$.
Here $b\ne0$ follows from the degree, and $a=0$ would give a factor
$XY$. If $a/b$ were real, $P$ would factor as
$(XY+a/b)(bX^m+\overline bY^m)$.
Thus $a/b\notin\R$. Substitute $z=r e^{i\phi}w$, choosing
$r^2=|a/b|$ and $e^{im\phi}b$ positive real, and rescale the equation
by a nonzero real constant. This gives~\eqref{eq:family} with
$\alpha=(a/b)/|a/b|$. The converse is Theorem~\ref{thm:family}.
\end{proof}

\section{Equality cases on the sphere}

Put
\begin{equation}\label{eq:complexification}
 P_\alpha(X,Y)=X^m(\alpha+XY)+Y^m(\overline\alpha+XY),
 \qquad m\ge2.
\end{equation}
Its closure $V_\alpha$ in $\PP^1\times\PP^1$ has bidegree
$(m+1,m+1)$. This compactification, rather than the ordinary projective
plane, lets ambient M\"obius maps act by regular automorphisms.

\begin{lemma}\label{lem:geometry}
For $|\alpha|=1$ and $\alpha\notin\R$, $P_\alpha$ is irreducible.
The only singular points of $V_\alpha$ are $(0,0)$ and
$(\infty,\infty)$, both ordinary $m$-fold points. Its normalization
has genus $m$, and $C_{m,\alpha}$ is infinite.
\end{lemma}

\begin{proof}
In the chart $X=tY$, remove the exceptional factor $Y^m$ to obtain
\begin{equation}\label{eq:doublecover}
 H(t,Y)=(\alpha t^m+\overline\alpha)+t(t^m+1)Y^2=0.
\end{equation}
The two coefficients in this quadratic in $Y$ are coprime in
$\C[t]$: at $t=0$ the first is nonzero, and at $t^m=-1$ it is
$\overline\alpha-\alpha\ne0$.
The rational function
\[
 -\frac{\alpha t^m+\overline\alpha}{t(t^m+1)}
\]
has a simple pole at $t=0$, so it is not a square in $\C(t)$.
Gauss's lemma makes $H$ irreducible. The change $X=tY$ is invertible
where $Y\ne0$, and $P_\alpha$ has no factor $Y$, proving its
irreducibility.

On $XY\ne0$, equation~\eqref{eq:doublecover} is a local coordinate
description. Its $Y$ derivative can vanish only at $t^m=-1$,
where $H\ne0$. Thus there are no singularities there.
The tangent cones at $(0,0)$ and $(\infty,\infty)$ are, respectively,
$\alpha X^m+\overline\alpha Y^m$ and $u^m+v^m$, with
$u=1/X$, $v=1/Y$. Each has $m$ distinct linear factors.
The other two boundary points are $(\infty,0)$ and $(0,\infty)$.
Near the first, the equation is
\[
 u^{m+1}P_\alpha(1/u,Y)
 =\alpha u+Y+\overline\alpha u^{m+1}Y^m+u^{m}Y^{m+1},
\]
whose linear part is nonzero. The other point is treated by exchanging
$X,Y$ and conjugating coefficients. Setting either coordinate to zero
or infinity shows that these are all boundary points.

The double cover~\eqref{eq:doublecover} has $2m+2$ simple branch
points: the $m$ roots of each of its coprime degree-$m$ factors, and
$0,\infty$. Riemann--Hurwitz gives $2g-2=-4+(2m+2)$, so $g=m$. % chktex 8
For each $r>0$ the equation
$\Rea(e^{im\theta}(r^2+\alpha))=0$ has $2m$ distinct angles
modulo $2\pi$, because $r^2+\alpha\ne0$.
\end{proof}

\begin{proof}[Proof of Theorem~\ref{thm:family}]
The degree and real-locus assertions follow from~\eqref{eq:complexification}
and Lemma~\ref{lem:geometry}. A holomorphic map $M$ acts on the
complexification by $(X,Y)\mapsto(M(X),\overline M(Y))$;
an anti-M\"obius map $M\circ\mathrm{conj}$ acts by
$(X,Y)\mapsto(M(Y),\overline M(X))$.
Here $\overline M$ conjugates the coefficients of $M$.
The infinite real locus is Zariski dense, so an equivalence of the
spherical real loci carries $V_\alpha$ onto $V_\beta$.
It must preserve the unordered pair of singular points. Consequently
the ambient map has one of the four forms
\[
 cz,\qquad c/z,\qquad c\overline z,\qquad c/\overline z
 \quad(c\in\C^\times).
\]

Write $t=|c|^2$. Substitute each form into $P_\beta$, clearing the
denominator by $X^{m+1}Y^{m+1}$ when needed. Proportionality to
$P_\alpha$ forces the ratio of the coefficients of $X^m$ and
$X^{m+1}Y$ to equal $\alpha$. In the four cases these ratios are
\begin{equation}\label{eq:ratios}
 \frac{\beta}{t},\qquad \frac{t}{\overline\beta},\qquad
 \frac{\overline\beta}{t},\qquad \frac{t}{\beta},
 \quad\text{respectively}.
\end{equation}
Since $|\alpha|=|\beta|=1$, necessarily $t=1$.
The holomorphic cases therefore require $\alpha=\beta$; the
antiholomorphic cases require $\alpha=\overline\beta$.
These conditions are sufficient: use the identity or conjugation.
In particular a self-equivalence cannot be antiholomorphic.

For $\beta=\alpha$, comparison of the remaining coefficients gives
$c^{2m}=1$ for $cz$, and $c^{2m}=\overline\alpha^{\,2}$ for $c/z$.
Substitution verifies both identities, proving~\eqref{eq:group}.
Choose $\zeta=e^{\pi i/m}$ and $c_0^m=\overline\alpha$.
The maps $r(z)=\zeta z$ and $s(z)=c_0/z$ satisfy
$r^{2m}=s^2=1$ and $srs=r^{-1}$; their $4m$ distinct maps are
exactly~\eqref{eq:group}. Among them only the rotations are Euclidean
isometries, since $c/z$ does not fix infinity.
\end{proof}

\begin{remark}[Scope of the classification]\label{rem:scope}
The restriction $d\ge5$ in the equality clause of
Theorem~\ref{thm:bounds} is necessary: in degree four,
$\Rea(z^4)=1$ also attains the rotation bound and is not in the
$m=2$ family, since its normalization has genus three rather than two.
For fixed $m\ge2$, the parameters $\alpha=e^{i\theta}$,
$0<\theta<\pi$, give pairwise fully M\"obius-inequivalent examples.
If $\alpha$ is algebraic, all the maps in~\eqref{eq:group} have
algebraic coefficients. These are explicit consequences of the
classification, not separate finiteness or descent theorems.
\end{remark}

\appendix
\section{Formal verification}\label{app:lean}

Every theorem, lemma and remark above is checked in Lean~4~\cite{Lean4} with
Mathlib~\cite{Mathlib}, in the repository
\url{https://github.com/carlok/curve-symmetry-lean}: 108 modules and about 900
theorems and lemmas, built with Lean~4.34.0 and Mathlib at commit
\texttt{5ed2965} (tag \texttt{v4.34.0}), with nine exact dependency pins. The
proofs use only the axioms \lean{propext}, \lean{Classical.choice} and
\lean{Quot.sound}, with no \lean{sorry} and no \lean{native_decide}; an audit
checks this for every declaration of the project. The statements start from
real Cartesian polynomials, the actual Euclidean isometry groups of their real
loci and actual M\"obius and anti-M\"obius maps of the Riemann sphere.

\paragraph{Where each statement is checked.}
\begin{itemize}[before=\raggedright]
\item Theorem~\ref{thm:bounds}: \lean{paper_sharp_bounds} (finiteness, the
  cyclic direct group and both bounds), \lean{paper_rotation_sharp} and
  \lean{paper_full_sharp} (sharpness), \lean{paper_equality_classification}
  and \lean{paper_family_converse} (the equality case).
\item Theorem~\ref{thm:family}: \lean{paper_family_geometry};
  \lean{family_mobius_complete} and \lean{family_anti_mobius_complete} (every
  equivalence has one of the four forms); \lean{family_mobius_self_mem_iff}
  and \lean{family_no_anti_mobius_self}; \lean{familyAmbientGroup_dihedral}
  and \lean{familyAmbientGroup_card}; \lean{family_isometry_card};
  \lean{family_mobius_equivalence_iff} and
  \lean{family_anti_mobius_equivalence_iff}.
\item Lemma~\ref{lem:sign}: \lean{paper_isometry_sign}. Equation~\eqref{eq:radial}:
  \lean{paper_high_order_radial_form}. The quintic~\eqref{eq:quintic}:
  \lean{quintic_isometry_count}.
\item Lemma~\ref{lem:geometry}: \lean{paper_family_geometry} (irreducibility),
  \lean{family_projective_critical_eq_pair} (the singular pair),
  \lean{family_ordinary_multiple_points}, \lean{paper_family_genus} and
  \lean{family_metric_circle_card}.
\item Remark~\ref{rem:scope}: \lean{paper_degree_four_genus},
  \lean{family_arc_equivalence_iff} and
  \lean{familyAmbientGroup_algebraic_representative}.
\end{itemize}
The file \texttt{COVERAGE.md} of the repository lists every claim, including
the supporting ones, with its declarations.

\paragraph{Readings.}
A few words of the text receive one precise meaning in Lean, and a reader
should check that it is the intended one. The normalization of $V_\alpha$ is
read as its function field, a double cover of $\C(t)$. Its genus is the
dimension of the holomorphic differentials of the function field, those
regular at every place, a place being a valuation ring containing $\C$. This
genus does not depend on a model, so an isomorphism of function fields
preserves it; a similarity between real loci induces such an isomorphism, and
this is how Remark~\ref{rem:scope} is checked, for orientation-preserving and
orientation-reversing similarities and every nonreal parameter. Branch points
are counted by fibres (one place over the point instead of two); ramification
indices are not formalized. Singular points, multiplicities and tangent cones
are computed in the four standard charts of $\PP^1\times\PP^1$, on complex
points with the classical Zariski topology; invariance under a change of
chart is not claimed.

\paragraph{Different routes.}
Lean obtains the genus in Lemma~\ref{lem:geometry} from an explicit basis
$t^i\,dt/w$, $0\le i<m$, of the holomorphic differentials, where
$w^2=-t(t^m+1)(\alpha t^m+\overline\alpha)$, instead of from Riemann--Hurwitz. % chktex 8
The genus three in Remark~\ref{rem:scope} comes the same way: in the
coordinates $X=z$, $Y=\overline z$ the curve $\Rea(z^4)=1$ becomes
$X^4+Y^4=2$, and $dX/Y^3$, $X\,dX/Y^3$ and $dX/Y^2$ form a basis. The
irreducibility of~\eqref{eq:doublecover} follows from Eisenstein's criterion at
$t=0$ after reversing the quadratic variable, in place of the nonsquare
argument.

\paragraph{What is checked, and by whom.}
The Lean kernel checks every statement listed above, and GitHub Actions
rebuilds the repository from a clean checkout. Theorem~\ref{thm:bounds} alone
is also published in \url{https://github.com/carlok/sharp-symmetry-bounds-lean}
and registered with the Palomar registry as
\texttt{PALOMAR-2026-09-18-000007}, after mechanical verification and an % chktex 8
automated editorial review; nothing else in this note is registered. The
readings above are the author's choices. No independent human review has
taken place, and neither the kernel checks nor the registration certify
novelty. The Lean code was written with AI coding agents under the author's
direction. With Elan installed, \texttt{lake exe cache get} and
\texttt{lake build} in the repository rebuild the proofs, and
\texttt{lake env lean verification/Audit.lean} reruns the axiom audit.

\begin{thebibliography}{9}
\raggedright % chktex 1
\bibitem{LRG}
P.~Lebmeir and J.~Richter-Gebert,
\emph{Rotations, translations and symmetry detection for complexified curves},
Computer Aided Geometric Design \textbf{25} (2008), 707--719.
\href{https://doi.org/10.1016/j.cagd.2008.09.004}{doi:10.1016/j.cagd.2008.09.004}.
Author manuscript:
\href{https://science-to-touch.com/Articles/jrg/53_VorauPaper.pdf}{Section~4}.

\bibitem{Lebmeir}
P.~M.~Lebmeir,
\emph{Feature Detection for Real Plane Algebraic Curves},
doctoral thesis, Technische Universit\"at M\"unchen, 2009.
\href{https://mediatum.ub.tum.de/doc/681056/512338.pdf}{Repository PDF}.

\bibitem{PZ}
J.~Pach and F.~de Zeeuw,
\emph{Distinct distances on algebraic curves in the plane},
\href{https://arxiv.org/abs/1308.0177v3}{arXiv:1308.0177v3} (2014),
Lemma~2.6.

\bibitem{ALVv1}
J.~G.~Alc\'azar, M.~L\'avi\v cka and J.~Vr\v sek,
\emph{Symmetries and similarities of planar algebraic curves using
harmonic polynomials},
\href{https://arxiv.org/abs/1801.09962v1}{arXiv:1801.09962v1} (2018),
Lemma~9. A journal article with this title appeared in
Journal of Computational and Applied Mathematics \textbf{357} (2019),
302--318, \href{https://doi.org/10.1016/j.cam.2019.02.036}{doi:10.1016/j.cam.2019.02.036};
the comparison here is with the arXiv version only.

\bibitem{Lean4}
L.~de~Moura and S.~Ullrich,
\emph{The Lean~4 theorem prover and programming language},
in Automated Deduction, CADE~28, Lecture Notes in Computer Science
\textbf{12699}, Springer, 2021, 625--635.
\href{https://doi.org/10.1007/978-3-030-79876-5_37}{doi:10.1007/978-3-030-79876-5\_37}.

\bibitem{Mathlib}
The mathlib Community,
\emph{The Lean mathematical library},
in Proceedings of the 9th ACM SIGPLAN International Conference on Certified
Programs and Proofs (CPP 2020), 367--381.
\href{https://doi.org/10.1145/3372885.3373824}{doi:10.1145/3372885.3373824}.
\end{thebibliography}
\end{document}
